%%%%%%%%%%%%%%%%%%%%%%%%%%%%%%%%%%%%%%%%%
% Jacobs Landscape Poster
% LaTeX Template
% Version 1.0 (29/03/13)
%
% Created by:
% Computational Physics and Biophysics Group, Jacobs University
% https://teamwork.jacobs-university.de:8443/confluence/display/CoPandBiG/LaTeX+Poster
% 
% Further modified by:
% Nathaniel Johnston (nathaniel@njohnston.ca)
%
% This template has been downloaded from:
% http://www.LaTeXTemplates.com
%
% License:
% CC BY-NC-SA 3.0 (http://creativecommons.org/licenses/by-nc-sa/3.0/)
%
%%%%%%%%%%%%%%%%%%%%%%%%%%%%%%%%%%%%%%%%%

%----------------------------------------------------------------------------------------
%	PACKAGES AND OTHER DOCUMENT CONFIGURATIONS
%----------------------------------------------------------------------------------------

\documentclass[final]{beamer}

\usepackage[scale=1.24]{beamerposter} % Use the beamerposter package for laying out the poster

\usetheme{confposter} % Use the confposter theme supplied with this template

\setbeamercolor{block title}{fg=ngreen,bg=white} % Colors of the block titles
\setbeamercolor{block body}{fg=black,bg=white} % Colors of the body of blocks
\setbeamercolor{block alerted title}{fg=white,bg=dblue!70} % Colors of the highlighted block titles
\setbeamercolor{block alerted body}{fg=black,bg=dblue!10} % Colors of the body of highlighted blocks
% Many more colors are available for use in beamerthemeconfposter.sty

%-----------------------------------------------------------
% Define the column widths and overall poster size
% To set effective sepwid, onecolwid and twocolwid values, first choose how many columns you want and how much separation you want between columns
% In this template, the separation width chosen is 0.024 of the paper width and a 4-column layout
% onecolwid should therefore be (1-(# of columns+1)*sepwid)/# of columns e.g. (1-(4+1)*0.024)/4 = 0.22
% Set twocolwid to be (2*onecolwid)+sepwid = 0.464
% Set threecolwid to be (3*onecolwid)+2*sepwid = 0.708

\newlength{\sepwid}
\newlength{\onecolwid}
\newlength{\twocolwid}
\newlength{\threecolwid}
\setlength{\paperwidth}{48in} % A0 width: 46.8in
\setlength{\paperheight}{36in} % A0 height: 33.1in
\setlength{\sepwid}{0.024\paperwidth} % Separation width (white space) between columns
\setlength{\onecolwid}{0.22\paperwidth} % Width of one column
\setlength{\twocolwid}{0.464\paperwidth} % Width of two columns
\setlength{\threecolwid}{0.708\paperwidth} % Width of three columns
\setlength{\topmargin}{-0.5in} % Reduce the top margin size
%-----------------------------------------------------------

\usepackage{graphicx}  % Required for including images

\usepackage{booktabs} % Top and bottom rules for tables

%----------------------------------------------------------------------------------------
%	TITLE SECTION 
%----------------------------------------------------------------------------------------

\title{Sandpiles} % Poster title

\author{Alex Pan} % Author(s)

\institute{Reed College} % Institution(s)

%----------------------------------------------------------------------------------------

\begin{document}

\addtobeamertemplate{block end}{}{\vspace*{2ex}} % White space under blocks
\addtobeamertemplate{block alerted end}{}{\vspace*{2ex}} % White space under highlighted (alert) blocks

\setlength{\belowcaptionskip}{2ex} % White space under figures
\setlength\belowdisplayshortskip{2ex} % White space under equations

\begin{frame}[t] % The whole poster is enclosed in one beamer frame

\begin{columns}[t] % The whole poster consists of three major columns, the second of which is split into two columns twice - the [t] option aligns each column's content to the top

\begin{column}{\sepwid}\end{column} % Empty spacer column

\begin{column}{\onecolwid} % The first column

%----------------------------------------------------------------------------------------
%	OBJECTIVES
%----------------------------------------------------------------------------------------

\begin{alertblock}{Abstract}

In the following, I will discuss the sandpile model and the closely-related phenomenon of self-organized criticality. I will discuss some features of sandpiles that we might want to know information about and the utility of simulation in obtaining measurements on sandpiles. I then present my own simulations of some 1D and 2D sandpile configurations and discuss the results. I will also dedicate some of the discussion to various issues I ran into when trying to get my simulations working.

\end{alertblock}

%----------------------------------------------------------------------------------------
%	INTRODUCTION
%----------------------------------------------------------------------------------------

\begin{block}{Introduction}

The \textit{sandpile} model was first introduced in 1978 by Bak, Tang, and Weisenfeld, as a system capable of displaying spatial and temporal power laws and scale-invariance without finely-tuned parameters. This behavior was described as \textit{self-organized criticality} (SOC) because of the way these systems can arrive at critical states on their own. SOC was given as the underlying process behind ubiquitous and complex phenomena like fractal geometry and ``$1/f$'' noise in the 1987 paper, and has since been used to model physical systems like earthquakes and rainfall, as well as non-physical systems in sociology and stock-markets.

\end{block}

%------------------------------------------------

\begin{block}{Criticality}
\textit{Critical points} exhibit scale-invariance through features like fractals and power laws, but cannot necessarily be reached depending on a system's parameters. For example, in the Ising model, the system doesn't display critical behavior unless the parameters have been carefully tuned. However, in systems that exhibit SOC, parameters can be changed widely without affecting the systems' ability to start exhibiting critical behavior.
\end{block}

%----------------------------------------------------------------------------------------

\end{column} % End of the first column

\begin{column}{\sepwid}\end{column} % Empty spacer column

\begin{column}{\onecolwid} % The first column

%----------------------------------------------------------------------------------------
%	SANDPILES
%----------------------------------------------------------------------------------------

\begin{block}{Sandpiles}

In very general terms, sandpiles can be thought of as cellular automota. There are three things that characterize the behavior of a specific configuration:
\begin{itemize}
\item A measure $z$ (not in the formal sense)
\item A critical value $K$
\item Some action $f$ by which a cell reduces its measure
\end{itemize} 
Whenever a cell's measure exceeds the critical value, it performs the action $f$. (i.e. when a cell $x$ has $z(x)>K$, $f(x)$ happens).

In the original paper, this is defined more strictly: when a cell's slope with its neighbors exceeds a critical value, it loses some of its weight to its nieghbors. However, I prefer to use looser defintions because I believe the stricter ones might give some false intuition about the measure. While slope is a good example because of its physical parallels, it suggests that the measure must be defined on a relationship between a cell and its neighbors. However this does not seem to be the case according to simulated results, as I will discuss later.

For the sake of simplicity and clarity, I focused on the 1D and 2D models used by Bak, Tang, and Weisenfeld, as well as some variants in the 2D case.

\end{block}

%----------------------------------------------------------------------------------------
%	SETUP
%----------------------------------------------------------------------------------------

\begin{block}{Setup}

As physical experiments have found, observing SOC on real sandpiles is not easy or necessarily possible. Instead, I use Monte Carlo simulations because they provide very high control over the parameters, which we expect to be important in observing SOC, and because they provide can provide fine-grain measurements that might not be easy to obtain in a physical setting.

In all configurations, I add grains to the system uniformly, because that seems to be the ``most'' random way of perturbing the system. The key measurements of interest are the size and length of avalanches. 


\end{block}

%----------------------------------------------------------------------------------------

\end{column} % End of column 2.2



%----------------------------------------------------------------------------------------


\begin{column}{\onecolwid} % The first column within column 2 (column 2.1)

%----------------------------------------------------------------------------------------
%	MATHEMATICAL SECTION
%----------------------------------------------------------------------------------------

\begin{block}{1D BTW Results}


\end{block}


%----------------------------------------------------------------------------------------
%	RESULTS
%----------------------------------------------------------------------------------------

\begin{block}{2D BTW Results}

% \begin{figure}
% \includegraphics[width=0.8\linewidth]{placeholder.jpg}
% \caption{Figure caption}
% \end{figure}

% Nunc tempus venenatis facilisis. Curabitur suscipit consequat eros non porttitor. Sed a massa dolor, id ornare enim:

% \begin{table}
% \vspace{2ex}
% \begin{tabular}{l l l}
% \toprule
% \textbf{Treatments} & \textbf{Response 1} & \textbf{Response 2}\\
% \midrule
% Treatment 1 & 0.0003262 & 0.562 \\
% Treatment 2 & 0.0015681 & 0.910 \\
% Treatment 3 & 0.0009271 & 0.296 \\
% \bottomrule
% \end{tabular}
% \caption{Table caption}
% \end{table}

\end{block}

%----------------------------------------------------------------------------------------

\end{column} % End of column 2.2


\begin{column}{\sepwid}\end{column} % Empty spacer column

\begin{column}{\onecolwid} % The third column


\begin{block}{Modified 2D Results}

\end{block}

%----------------------------------------------------------------------------------------
%	CONCLUSION
%----------------------------------------------------------------------------------------

\begin{block}{Conclusion}



\end{block}


%----------------------------------------------------------------------------------------
%	REFERENCES
%----------------------------------------------------------------------------------------

\begin{block}{References}

\nocite{*} % Insert publications even if they are not cited in the poster
\small{\bibliographystyle{unsrt}
\bibliography{sample}\vspace{0.75in}}

\end{block}


%----------------------------------------------------------------------------------------
%	CONTACT INFORMATION
%----------------------------------------------------------------------------------------

\setbeamercolor{block alerted title}{fg=black,bg=norange} % Change the alert block title colors
\setbeamercolor{block alerted body}{fg=black,bg=white} % Change the alert block body colors

\begin{alertblock}{Simulations}

\begin{itemize}
\item Github: \href{https://github.com/apanana/Sandpiles}{https://github.com/apanana/Sandpiles}
\end{itemize}

There are several Mathematica notebooks containing both data and user-interactive dynamic modules that allow for further manipulation of parameters
\end{alertblock}

%----------------------------------------------------------------------------------------

\end{column} % End of the third column

\end{columns} % End of all the columns in the poster

\end{frame} % End of the enclosing frame

\end{document}
